A finite graph $F$ is Ramsey size-linear if $r(F,H)=O_F(e(H))$
for every nonempty graph $H$ with no isolated vertices. Are there
infinitely many graphs which fail this property although every
proper subgraph has it? The answer is yes. We reconstruct Wigderson's
argument and prove the elementary forest and large-girth ingredients.
Graphs are considered up to isomorphism throughout.

The external Ramsey lower bound used is the following standard
consequence of the Lov\'asz local lemma, stated in Wigderson's note:
for a fixed graph $F$ with $v=v(F)\geq3$ and $e=e(F)$,
\[
 r(F,K_q)=\Omega_F\left((q/\log q)^{(e-1)/(v-2)}\right).
\]
Only graphs satisfying $e\geq2v-2$ are needed. For these,
$(e-1)/(v-2)\geq2+1/(v-2)>2$, and hence
$r(F,K_q)/e(K_q)\to\infty$. Such $F$ are not Ramsey size-linear.
The cited lower bound, rather than its probabilistic proof, is
the sole non-elementary dependency of this reconstruction.
